\section{工程视角：为什么 PPO 成为默认选项}
\subsection{复杂度、稳定性与可实现性的折中}
课程虽然只用 15 分钟讲解 PG、TRPO、PPO 的演化，但背后的工程结论非常明确：PPO 之所以流行，不是因为它在所有 benchmark 上都绝对最优，而是因为它在\textbf{稳定性、实现门槛和训练效率}之间取得了最平衡的折中。

\begin{table}[H]
\centering
\begin{tabular}{p{2.6cm}p{3.5cm}p{3.7cm}p{3.2cm}}
\toprule
方法 & 优点 & 主要代价 & 适用判断 \\
\midrule
PG & 概念最直接 & 方差高、更新不稳 & 用于理解基本原理 \\
TRPO & 理论约束更强 & 二阶近似和实现复杂 & 更适合研究型推导 \\
PPO & 稳定且实现简单 & 仍需精调 clip/advantage & 工程实践的默认起点 \\
\bottomrule
\end{tabular}
\caption{PG、TRPO、PPO 在工程中的典型取舍}
\end{table}

\begin{warningbox}{不要把 PPO 当作“免调参算法”}
PPO 只是把 TRPO 的复杂约束换成了更易用的剪裁目标，并不意味着训练就会自动稳定。batch 大小、优势归一化、reward 尺度和 rollout 质量依旧会直接影响结果。
\end{warningbox}

\subsection{本章小结}
从工程落地角度看，PPO 是一个高性价比基线：先用它搭通训练闭环，再决定是否有必要引入更复杂的约束或更新规则，通常比一开始就追求“最强算法”更现实。

\newpage
